%% toc: P3 25 25.1
\section{Vaught's Results on Theories Categorical in an Uncountable Power}
\label{sec:p03-w03-p03-007-vaught-categorical}

The fifth section of Vaught's 1961 paper applies the prime-model and
saturation criteria of its first four sections to theories with few models
of some infinite power \cite[\S~5, pp.~315--318]{Vaught1961}. It takes as
input a type-counting theorem of Ehrenfeucht. Its conclusions are that such
theories have prime models and countable universal homogeneous models; that,
for a theory categorical in \(\aleph_1\) but not in \(\aleph_0\), the prime
model has no proper elementary submodel; and that, in a theory categorical
in \(\aleph_1\), a parameter-free definable set of one variable is either
finite or as large as the model. The same section illustrates the first
conclusion with algebraically closed fields. The two examples that Vaught
gives just before the section, discrete orders and real closed fields, are
recorded in \cref{sec:p03-w03-p03-007-illustrations} with the standing the
source gives them.

This section states those results with the hypotheses and cardinal scopes
the source uses and proves them in the annex from the constructions of
\cref{sec:p03-n-a-003-01}. Ehrenfeucht's theorem is stated with its
attribution and used as a cited result. For theories that are categorical
in some infinite power, the conclusion of that theorem is also derived here
from results proved elsewhere in this part, so that the categorical cases
of Theorems~5.3 and~5.4 do not depend on the cited proof. Vaught's own
remarks on what he could not prove in 1961 are recorded with their status.

\subsection{Setting, notation and the cited type-counting theorem}
\label{sec:p03-w03-p03-007-setting}

\begin{hbconvention}{def:p03-w03-p03-007-setting}{Standing theory, powers and definable sets}
The theory \(T\) is as in \cref{def:p03-vaught-standing}: complete, in an at
most countable language \(L\) of relation symbols and individual constants,
with identity, and with an infinite model, so that every model of \(T\) is
infinite. The formula sets \(F_n(T)\) and the complete-type sets
\(\mathcal P_n(T)\) are those of \cref{def:p03-vaught-formula-algebra,def:p03-vaught-prime-ideals};
prime, atomic, \(\aleph_0\)-universal, homogeneous, weakly saturated and
saturated models are as in
\cref{def:p03-vaught-model-notions,def:p03-vaught-saturation}.
The \emph{power} of a model is the cardinality of its universe.
For an infinite cardinal \(\kappa\), \(T\) is \emph{\(\kappa\)-categorical}
when any two models of \(T\) of power \(\kappa\) are isomorphic. The source
writes \(\aleph_\alpha\) for an arbitrary infinite power and
\(\aleph_{1+\alpha}\) for an arbitrary uncountable one.

For \(\theta\in F_1(T)\) and \(\mathcal A\models T\), put
\[
 \theta(\mathcal A)=\{x\in\dom\mathcal A:\mathcal A\models\theta[x]\},
\]
the source's \(\{x/{\models_{\mathcal A}}\theta[x]\}\). The formula
\(\theta\) has no parameters. A model is \emph{minimal} when it has no
proper elementary submodel, that is, no \(\mathcal A'\preccurlyeq\mathcal A\)
with \(\dom\mathcal A'\neq\dom\mathcal A\).
The construction of \cref{sec:p03-w03-p03-007-definable} uses the ambient
Axiom of Choice of \cref{def:p03-vaught-standing}, and so do the cardinal
comparisons below.
\end{hbconvention}

A note to the title of Vaught's fifth section points to examples of
categorical theories in a paper of Łoś (1955) and in Vaught's own 1954
paper, and repeats their observation that such a theory without finite
models must be complete \cite[p.~315, n.~17; p.~321, refs.~9 and~19]{Vaught1961}.
For a countable language the observation reads as follows; the standing
theory is complete by assumption, so the lemma is recorded for theories
that are not assumed complete.

\begin{hblemma}{thm:p03-w03-p03-007-categorical-complete}{proved=annex:p03-w03-p03-007-categorical-complete}{Categoricity without finite models gives completeness (Vaught's note~17)}
Let \(T_0\) be a set of sentences of an at most countable language \(L\)
that has a model and has no finite model. If \(T_0\) is
\(\kappa\)-categorical for some infinite cardinal \(\kappa\), then \(T_0\) is
complete: for every sentence \(\sigma\) of \(L\), every model of \(T_0\)
satisfies \(\sigma\) or every model of \(T_0\) satisfies \(\neg\sigma\).
\end{hblemma}

Vaught opens his fifth section with the following theorem, which he
credits to Ehrenfeucht \cite[p.~315, 5.1]{Vaught1961}.

\begin{hbtheorem}{thm:p03-w03-p03-007-ehrenfeucht}{cited=Ehrenfeucht1958Notices}{Ehrenfeucht's theorem (Vaught's 5.1)}
If, for some ordinal \(\alpha\), the theory \(T\) has fewer than
\(2^{\aleph_0}\) pairwise non-isomorphic models of power \(\aleph_\alpha\),
then every \(\mathcal P_n(T)\), \(n<\omega\), is countable.
\end{hbtheorem}

The attribution and standing of this theorem are those recorded in Vaught's
footnote~18 \cite[p.~315, n.~18]{Vaught1961}. Vaught refers the result to
an abstract of Ehrenfeucht \cite{Ehrenfeucht1958Notices}, where only the
case \(n=1\) is stated. He adds that an earlier paper of Ehrenfeucht had
shown that a theory categorical in power \(2^{\aleph_\alpha}\) has an
\(\aleph_0\)-universal model. The case \(\alpha=0\) is, in Vaught's words,
``much more easily proved''; he reports that Mostowski had established it for
some special theories by a method that works for any \(T\), and he recalls
that a countable Boolean algebra has either countably many or
\(2^{\aleph_0}\) prime ideals. Neither Vaught nor this section supplies a
proof of the general theorem; every use of it below is marked. In the
categorical case the conclusion is obtained independently as follows.

\begin{hblemma}{thm:p03-w03-p03-007-categorical-small}{proved=annex:p03-w03-p03-007-categorical-small}{Categoricity in one infinite power gives countable type sets}
If \(T\) is \(\kappa\)-categorical for some infinite cardinal \(\kappa\),
then every \(\mathcal P_n(T)\) is countable; that is, \(T\) is small in the
sense of \cref{def:p03-vaught-smallness}. For \(\kappa=\aleph_0\) every
\(\mathcal P_n(T)\) is in fact finite.
\end{hblemma}

A \(\kappa\)-categorical theory satisfies the hypothesis of
Ehrenfeucht's theorem with \(\aleph_\alpha=\kappa\), since it has exactly one
model of that power up to isomorphism and \(1<2^{\aleph_0}\)
\cite[p.~316]{Vaught1961}. The lemma reaches the same conclusion for these
theories without the cited proof. It does not cover the general hypothesis
of fewer than \(2^{\aleph_0}\) models, which only the cited theorem treats.

\subsection{Prime and universal homogeneous models (Vaught's 5.2)}
\label{sec:p03-w03-p03-007-five-two}

\begin{hbtheorem}{thm:p03-w03-p03-007-five-two}{proved=annex:p03-w03-p03-007-five-two}{Prime and universal homogeneous models under few models (Vaught's 5.2)}
If the hypothesis of \cref{thm:p03-w03-p03-007-ehrenfeucht} holds, then
\(T\) has an \(\aleph_0\)-universal, homogeneous model and a prime model.
Both are countably infinite and each is unique up to isomorphism. If \(T\) is
\(\kappa\)-categorical for some infinite \(\kappa\), the same conclusion holds
with \cref{thm:p03-w03-p03-007-categorical-small} in place of the cited
theorem.
\end{hbtheorem}

Vaught observes that once 5.1 gives countable type sets, the results of his
third and fourth sections apply, and he presents 5.2 as an immediate
consequence of his 3.5, 4.8 and 5.1 \cite[p.~315]{Vaught1961}. In this part
those are \cref{thm:p03-vaught-isolation-density,thm:p03-vaught-universal-prime};
the annex proof also uses his 4.4, 4.6 and 4.7
(\cref{thm:p03-vaught-saturated-uniqueness,thm:p03-vaught-saturation-characterization,thm:p03-vaught-smallness-equivalences})
for the existence and uniqueness of the universal homogeneous model.
The universal homogeneous model here is the countably infinite saturated
model of \cref{thm:p03-vaught-smallness-equivalences}, and the prime model
is the countably infinite atomic model of \cref{thm:p03-vaught-prime-atomic}.
The theorem does not say that they coincide. When \(T\) is
\(\aleph_0\)-categorical they do coincide, since \(T\) then has only one
countable model, and Vaught notes that the theorem has no interest in that
case \cite[p.~316]{Vaught1961}. His example of a theory categorical in the
uncountable powers but not in \(\aleph_0\) is the theory of algebraically
closed fields, treated in \cref{sec:p03-w03-p03-007-acf}.

\subsection{Elementary chains, copies and the minimality of prime models}
\label{sec:p03-w03-p03-007-minimal}

The proofs of Theorems~5.3 and~5.4 build chains of length \(\omega_1\)
whose countable stages keep a property of the first model. The two lemmas
below record the chain facts the source quotes from its Lemma~1.1, and the
renaming step that lets one extension pattern be repeated at every stage.

\begin{hblemma}{thm:p03-w03-p03-007-chains}{proved=annex:p03-w03-p03-007-chains}{Continuous elementary chains of models of \(T\)}
Let \(\gamma\) be an ordinal and let \(\langle\mathcal A_\xi:\xi<\gamma\rangle\)
be models of \(T\) with \(\mathcal A_\xi\preccurlyeq\mathcal A_{\xi+1}\)
whenever \(\xi+1<\gamma\), and with
\(\mathcal A_\eta=\bigcup_{\xi<\eta}\mathcal A_\xi\) for every limit
\(\eta<\gamma\).
\begin{enumerate}
\item \(\mathcal A_\xi\preccurlyeq\mathcal A_\zeta\) whenever
      \(\xi<\zeta<\gamma\).
\end{enumerate}
Suppose next that \(\delta\) is a limit ordinal, that
\(\mathcal A_\xi\preccurlyeq\mathcal A_\zeta\) for \(\xi<\zeta<\delta\),
and let \(\mathcal A=\bigcup_{\xi<\delta}\mathcal A_\xi\).
\begin{enumerate}\setcounter{enumi}{1}
\item Each \(\mathcal A_\xi\preccurlyeq\mathcal A\), and \(\mathcal A\models T\).
\item If every \(\mathcal A_\xi\) is atomic, so is \(\mathcal A\).
\item If every \(\mathcal A_\xi\) is saturated, so is \(\mathcal A\).
\item If \(\delta\) is countable and every \(\mathcal A_\xi\) is countable,
      \(\mathcal A\) is countably infinite.
\item For every \(\theta\in F_1(T)\),
      \(\theta(\mathcal A)=\bigcup_{\xi<\delta}\theta(\mathcal A_\xi)\).
\item If \(\delta=\omega_1\), every \(\mathcal A_\xi\) is countable, and
      \(\dom\mathcal A_{\xi+1}\neq\dom\mathcal A_\xi\) for every \(\xi\),
      then \(\mathcal A\) has power \(\aleph_1\).
\end{enumerate}
\end{hblemma}

\begin{hblemma}{thm:p03-w03-p03-007-copy}{proved=annex:p03-w03-p03-007-copy}{Copying an elementary extension over an isomorphic model}
Let \(\mathcal D\preccurlyeq\mathcal C\) be \(L\)-structures and let
\(g\colon\mathcal D\to\mathcal D'\) be an isomorphism. There are an
\(L\)-structure \(\mathcal C'\) and an isomorphism
\(F\colon\mathcal C\to\mathcal C'\) extending \(g\) such that
\(\mathcal D'\preccurlyeq\mathcal C'\) and
\(\dom\mathcal C'\setminus\dom\mathcal D'=F[\dom\mathcal C\setminus\dom\mathcal D]\).
In particular \(\mathcal D'\) is a proper elementary submodel of
\(\mathcal C'\) exactly when \(\mathcal D\) is one of \(\mathcal C\), and
\(\theta(\mathcal C')=F[\theta(\mathcal C)]\) for every \(\theta\in F_1(T)\).
\end{hblemma}

\begin{hbtheorem}{thm:p03-w03-p03-007-minimal-prime}{proved=annex:p03-w03-p03-007-minimal-prime}{Prime models are minimal (Vaught's 5.3)}
Suppose that \(T\) is \(\aleph_1\)-categorical but not
\(\aleph_0\)-categorical. Then every prime model of \(T\) is minimal: it has
no proper elementary submodel.
\end{hbtheorem}

A prime model exists under these hypotheses by
\cref{thm:p03-w03-p03-007-five-two}. The proof supposes a proper elementary
submodel, uses it to build an elementary chain of prime models of length
\(\omega_1\), and compares the atomic union of power \(\aleph_1\) with a model
of the same power that realizes a nonprincipal type. Both hypotheses are
used: \(\aleph_1\)-categoricity to identify the two models, and the failure of
\(\aleph_0\)-categoricity to supply the nonprincipal type through
\cref{thm:p03-vaught-finite-formulas-categoricity}. The second hypothesis
cannot be dropped (\cref{thm:p03-w03-p03-007-pure-sets}).

\begin{hbwarning}{def:p03-w03-p03-007-los-remark}{The cardinal \(\aleph_1\) in Vaught's statements}
Vaught's Theorems~5.3 and~5.4 assume categoricity in \(\aleph_1\)
specifically. He records the conjecture of Łoś that a theory categorical in
one non-denumerable power is categorical in all of them, observes that it
would allow \(\aleph_1\) to be replaced by \(\aleph_{1+\alpha}\), and states
that he could not prove that stronger version \cite[p.~316]{Vaught1961}.
The printed sentence names Theorem~5.2 at this point; Theorem~5.2 contains
no \(\aleph_1\), and the sentence that follows speaks of the stronger version
of 5.3, so the remark concerns 5.3. For 5.4 he again reports that he cannot
replace \(\aleph_1\) by \(\aleph_{1+\alpha}\), and states without giving the
argument that a similar but simpler proof allows \(\aleph_1\)-categoricity to
be replaced by categoricity in both \(\aleph_\alpha\) and \(\aleph_{\alpha+1}\)
\cite[p.~318, n.~22]{Vaught1961}. Vaught does not print that argument. The
claim of the footnote is nevertheless covered here: since
\(\aleph_{\alpha+1}\) is uncountable, categoricity in \(\aleph_{\alpha+1}\)
alone already gives the conclusion of 5.4 by
\cref{thm:p03-w03-p03-007-morley-transfer}. That corollary also gives the
stronger versions of 5.3 and 5.4 for every uncountable power. It rests on
Morley's categoricity theorem and is not one of Vaught's 1961 results.
\end{hbwarning}

\subsection{Definable sets in \(\aleph_1\)-categorical theories (Vaught's 5.4)}
\label{sec:p03-w03-p03-007-definable}

Vaught proves 5.4 by placing two saturated models inside one countable model
of an auxiliary theory. A new unary predicate \(V\) names the universe of an
elementary submodel, and the theory says that this submodel is proper and
contains every element satisfying \(\theta\). He writes the proof for a
single ternary relation symbol and says the general case is the same; the
definition below treats an arbitrary countable relational/constant
language directly.

\begin{hbdefinition}{def:p03-w03-p03-007-auxiliary-theory}{Relativization to \(V\) and the theory \(T'\)}
Let \(\theta\in F_1(T)\) and let \(V\) be a unary relation symbol not in
\(L\); put \(L_V=L\cup\{V\}\). For an \(L\)-formula \(\varphi\), its
relativization \(\varphi^V\) is obtained by leaving atomic formulas
unchanged, commuting with the connectives, and replacing
\(\exists v_j\psi\) by \(\exists v_j(Vv_j\land\psi^V)\) and
\(\forall v_j\psi\) by \(\forall v_j(Vv_j\to\psi^V)\).
This is relativization of the quantifiers to \(V\) in the sense
the source takes from Tarski, Mostowski and Robinson
\cite[p.~317, n.~21]{Vaught1961}.

The theory \(T'\) in \(L_V\) consists of the sentences of \(T\), the three
sentences
\[
 \exists v_0\,Vv_0,\qquad
 \exists v_1\,\neg Vv_1,\qquad
 \forall v_0\,[\theta(v_0)\to Vv_0],
 \tag{1}
\]
and, for every \(n<\omega\) and every \(\varphi\in F_n(T)\), the sentence
\[
 \forall v_0\cdots\forall v_{n-1}\,
 [Vv_0\land\cdots\land Vv_{n-1}\to(\varphi\leftrightarrow\varphi^V)],
 \tag{2}
\]
which for \(n=0\) is \(\varphi\leftrightarrow\varphi^V\).
The two relative interpretations of \(T\) in \(T'\), in the sense of
\cref{def:p03-vaught-relative-interpretation}, are \(I_=\), with domain
formula \(v_0=v_0\), and \(I_V\), with domain formula \(Vv_0\); both
interpret every relation symbol and every individual constant of \(L\) by
itself. For a model \(\mathfrak B\) of \(T'\) with \(L\)-reduct
\(\mathcal B\), \(\mathcal B_V\) denotes the \(L\)-structure with universe
\(V^{\mathfrak B}\) and the restricted relations.
\end{hbdefinition}

\begin{hblemma}{thm:p03-w03-p03-007-auxiliary}{proved=annex:p03-w03-p03-007-auxiliary}{What the auxiliary theory records}
Let \(\theta\), \(V\) and \(T'\) be as in
\cref{def:p03-w03-p03-007-auxiliary-theory}.
\begin{enumerate}
\item If \(\mathcal C\) is an \(L\)-structure, \(W\subseteq\dom\mathcal C\)
      is the universe of a substructure \(\mathcal C_W\), and
      \(\bar a\in W^n\), then
      \(\mathcal C_W\models\varphi[\bar a]\) if and only if
      \((\mathcal C,W)\models\varphi^V[\bar a]\), for every
      \(\varphi\in F_n(T)\).
\item Suppose that \(\mathcal C\models T\), that
      \(\mathcal D\preccurlyeq\mathcal C\) with
      \(\dom\mathcal D\neq\dom\mathcal C\), and that
      \(\theta(\mathcal C)\subseteq\dom\mathcal D\). Then
      \((\mathcal C,\dom\mathcal D)\models T'\).
\item If \(\mathfrak B=(\mathcal B,V^{\mathfrak B})\models T'\), then every
      individual constant of \(L\) denotes an element of \(V^{\mathfrak B}\),
      \(\mathcal B_V\) is a proper elementary submodel of \(\mathcal B\),
      both are models of \(T\), and
      \(\theta(\mathcal B)=\theta(\mathcal B_V)\).
\item \(I_=\) and \(I_V\) are relative interpretations of \(T\) in \(T'\),
      with \(I_=(\mathfrak B)=\mathcal B\) and
      \(I_V(\mathfrak B)=\mathcal B_V\) for every \(\mathfrak B\models T'\).
      The language \(L_V\) is countable.
\end{enumerate}
\end{hblemma}

\begin{hbtheorem}{thm:p03-w03-p03-007-definable-sets}{proved=annex:p03-w03-p03-007-definable-sets}{Definable sets are finite or of full power (Vaught's 5.4)}
Suppose that \(T\) is \(\aleph_1\)-categorical and \(\theta\in F_1(T)\).
Then, in every model \(\mathcal A\) of \(T\), the set \(\theta(\mathcal A)\)
is either finite or of the same power as \(\mathcal A\).
\end{hbtheorem}

The proof runs as follows. If a model had an infinite \(\theta\)-set of
smaller power, the auxiliary theory \(T'\) would be consistent. Since \(T\) is
small, Vaught's Theorem~4.2 then gives a countable model of \(T'\) in which
both interpretations are saturated. By the uniqueness of countable saturated
models, every countable saturated model of \(T\) therefore has a proper
countable saturated elementary extension with the same \(\theta\)-set.
Repeating this \(\omega_1\) times gives a model of power \(\aleph_1\) whose
\(\theta\)-set is countable. Compactness gives another model of power
\(\aleph_1\) whose \(\theta\)-set has power \(\aleph_1\), and the two models
contradict \(\aleph_1\)-categoricity. The source obtains smallness from
Ehrenfeucht's theorem \cite[p.~317]{Vaught1961}; here it comes from
\cref{thm:p03-w03-p03-007-categorical-small}. Vaught presents 5.4 as a
consequence that would follow easily from the conjecture of Łoś and that
might help to establish it \cite[p.~316]{Vaught1961}.

\begin{hbcorollary}{thm:p03-w03-p03-007-morley-transfer}{proved=annex:p03-w03-p03-007-morley-transfer}{Arbitrary uncountable powers}
Suppose that \(T\) is \(\kappa\)-categorical for some uncountable \(\kappa\).
Then the conclusion of \cref{thm:p03-w03-p03-007-definable-sets} holds
for \(T\), and if \(T\) is not \(\aleph_0\)-categorical, the conclusion of
\cref{thm:p03-w03-p03-007-minimal-prime} holds for \(T\).
\end{hbcorollary}

This corollary is a consequence of Morley's theorem
(\cref{thm:p03-morley-original-categoricity}), which resolves the conjecture
of Łoś for theories in a countable language. It is recorded here as a
relationship and is not part of Vaught's paper.

\subsection{The algebraically closed field illustration}
\label{sec:p03-w03-p03-007-acf}

Vaught illustrates 5.2 with ``the theory of algebraically closed fields'',
noting that it is categorical in the non-denumerable powers but not in
\(\aleph_0\). In it, he says, the field of complex algebraic numbers is a
prime model, and the algebraically closed fields of transcendence degree
\(\aleph_0\) are \(\aleph_0\)-universal homogeneous models
\cite[p.~316]{Vaught1961}. For these assertions he refers to Robinson's
\emph{Complete Theories} and to Tarski and Vaught's 1957 paper. A complete
theory needs the characteristic to be fixed, and the complex algebraic
numbers have characteristic \(0\). The theorem below therefore states and
proves the illustration for \(\mathsf{ACF}_0\). Its language
\(\{0,1,+,-,\cdot\}\) has function symbols. Vaught's results are stated for
relations and constants, and \cref{thm:p03-vaught-functional-bridge}
transfers them, together with primeness, saturation and homogeneity, to this
language.

\begin{hbtheorem}{thm:p03-w03-p03-007-acf}{proved=annex:p03-w03-p03-007-acf}{Prime and universal homogeneous algebraically closed fields}
Let \(\mathsf{ACF}_0\) be the theory of algebraically closed fields of
characteristic \(0\).
\begin{enumerate}
\item \(\mathsf{ACF}_0\) is complete, has only infinite models, is
      categorical in every uncountable power, and is not
      \(\aleph_0\)-categorical.
\item For \(K\models\mathsf{ACF}_0\), the set \(\mathbb A_K\) of elements of
      \(K\) algebraic over the prime field \(\mathbb Q\) is a countable
      algebraically closed subfield. It is a prime model of
      \(\mathsf{ACF}_0\), and it is minimal. Any two such fields are
      isomorphic. For \(K=\mathbb C\) it is the field of complex algebraic
      numbers.
\item A countable \(K\models\mathsf{ACF}_0\) has transcendence degree
      \(\aleph_0\) over \(\mathbb Q\) if and only if it is saturated, if and
      only if it is \(\aleph_0\)-universal and homogeneous. Such a field exists
      and is unique up to isomorphism.
\end{enumerate}
\end{hbtheorem}

The third clause includes a fact that Vaught does not state: the countable
fields of finite transcendence degree, the prime model among them, are not
\(\aleph_0\)-universal. The prime model is minimal, as 5.3 predicts, and this
can also be seen directly: an elementary submodel contains every element
algebraic over \(\mathbb Q\). In the case \(K=\mathbb C\), \(\mathbb C\) is a
model of \(\mathsf{ACF}_0\) by the fundamental theorem of algebra.

\subsection{Vaught's illustrations by discrete orders and real closed fields}
\label{sec:p03-w03-p03-007-illustrations}

Immediately before his fifth section, Vaught gives two theories as
illustrations of the hypotheses of his 3.5 and 4.7, the existence criteria
for prime models and for countable saturated models
(\cref{thm:p03-vaught-isolation-density,thm:p03-vaught-smallness-equivalences}).
He adds that his third and fourth sections give no new information about
either theory \cite[p.~315]{Vaught1961}. Every claim in the two examples is
asserted in the source without proof, with pointers to other work. The
record below keeps each claim in that standing and separates the one that is
proved in this handbook.

\begin{proofboundary}{def:p03-w03-p03-007-illustrations}{Vaught's two illustrations on page 315}
The following are Vaught's assertions \cite[p.~315 and n.~16; p.~321,
refs.~12, 16 and~18]{Vaught1961}, in his notation.
\begin{enumerate}
\item \emph{The theory \(T_1\).} \(T_1\) is the theory of infinite,
      discretely ordered systems that have a first element and no last
      element. Vaught introduces this choice of endpoints with ``say'', so
      the example is one instance of the discrete orders he has in mind.
\item \emph{Countable type sets for \(T_1\).} Each \(\mathcal P_n(T_1)\) is
      countable. Vaught calls this easily verified: the known decision
      procedure for \(T_1\) describes all possible definable sets and
      relations in models of \(T_1\). For a brief discussion of \(T_1\), and
      for references, his note~16 points to pp.~90--91 of Tarski and
      Vaught's 1957 paper \emph{Arithmetical extensions of relational
      systems}.
\item \emph{The universal homogeneous model of \(T_1\).} The
      \(\aleph_0\)-universal, homogeneous model of \(T_1\) is the system of
      order type \(\omega+(\omega^*+\omega)\cdot\eta\). Here \(\omega^*\) is
      the reverse of \(\omega\) and \(\eta\) is the order type of the
      rational numbers; the source uses these order-type symbols without
      defining them.
\item \emph{Many denumerable models of \(T_1\).} Vaught remarks that
      \(T_1\) has \(2^{\aleph_0}\) non-isomorphic denumerable models.
\item \emph{The theory \(T_2\) and its one-types.} \(T_2\) is the theory of
      real closed fields. The set \(\mathcal P_1(T_2)\) has \(2^{\aleph_0}\)
      members, which Vaught says is obvious.
\item \emph{Prime model and atomisticity for \(T_2\).} The field of real
      algebraic numbers is a prime model of \(T_2\), and each
      \(F_n(T_2)\) is atomistic. Vaught introduces both with ``as is
      known'' and cites Tarski's \emph{A decision method for elementary
      algebra and geometry} (second edition, 1951) and Robinson's
      \emph{Complete Theories} (1956).
\end{enumerate}
Vaught's results assume a complete theory, and page~315 does not argue
that \(T_1\) or \(T_2\) is complete. Items~2, 3, 4 and~6, and the
completeness of \(T_1\) and \(T_2\), are recorded as Vaught's attributed
claims. They are not proved in the source, and they are not proved in this
section. Item~5 is proved below for the complete theory of the real field;
its identification with \(T_2\) is the completeness of the theory of real
closed fields, one of the attributed claims.
\end{proofboundary}

\begin{hblemma}{thm:p03-w03-p03-007-real-types}{proved=annex:p03-w03-p03-007-real-types}{The real field has continuum many one-types}
Let \(T_{\mathbb R}\) be the complete theory of the field
\((\mathbb R;0,1,+,-,\cdot)\). Distinct real numbers realize distinct
complete one-types of \(T_{\mathbb R}\) over the empty set, and
\(T_{\mathbb R}\) has exactly \(2^{\aleph_0}\) complete one-types over the
empty set. Through \cref{thm:p03-vaught-functional-bridge}, the set
\(\mathcal P_1\) of the graph-language presentation of \(T_{\mathbb R}\)
has \(2^{\aleph_0}\) members.
\end{hblemma}

The examples are paired because they fall on opposite sides of the gap
between 3.5 and 4.7. If \(T_2\) is complete, the first part of item~6
implies the second by \cref{thm:p03-vaught-isolation-density}, and either
part says that \(T_2\) meets the condition of 3.5. By item~5 and
\cref{thm:p03-vaught-smallness-equivalences}, \(T_2\) has no denumerable
saturated model and no \(\aleph_0\)-universal model, so it fails the
hypothesis of 4.7. On Vaught's claims, then, a prime model can exist when
the one-type set is uncountable. For \(T_1\), item~2 says that \(T_1\) is
small, the first condition of 4.7, so
\cref{thm:p03-vaught-smallness-equivalences} yields a denumerable saturated
model, unique up to isomorphism by \cref{thm:p03-vaught-saturated-uniqueness}
and homogeneous by \cref{thm:p03-vaught-saturated-homogeneity}; the
identification of its order type in item~3 rests on the decision procedure.
Items~2 and~4 together would show that countable type sets do not bound
the number of denumerable models. Finite type sets do, by
\cref{thm:p03-vaught-finite-formulas-categoricity}: they hold exactly when
all denumerable models are isomorphic. Vaught's sixth section
(\cref{sec:p03-vaught-never-two-presentation}) shows that the number of
denumerable models of a complete theory is never exactly two.

The structure \((\omega,<)\) is a model of \(T_1\). For its complete theory
\(\operatorname{Th}(\omega,<)\), \cref{thm:p03-vaught-discrete-example}
proves directly that \((\omega,<)\) is atomic and prime, without the
decision procedure. That \(\operatorname{Th}(\omega,<)\) coincides with the
theory \(T_1\) given by Vaught's description is the completeness of \(T_1\),
which is not proved here.

\subsection{Boundary cases}
\label{sec:p03-w03-p03-007-boundaries}

\begin{hbtheorem}{thm:p03-w03-p03-007-pure-sets}{proved=annex:p03-w03-p03-007-pure-sets}{An \(\aleph_0\)-categorical theory with a non-minimal prime model}
Let \(T_\infty\) be the complete theory of infinite sets in the language
with identity only. It is categorical in every infinite power. Its prime
models are the countably infinite sets, and \(\omega\setminus\{0\}\) is a
proper elementary submodel of the prime model \(\omega\). Hence the
hypothesis ``not \(\aleph_0\)-categorical'' cannot be removed from
\cref{thm:p03-w03-p03-007-minimal-prime}.
\end{hbtheorem}

\begin{hbtheorem}{thm:p03-w03-p03-007-two-cardinal-boundary}{proved=annex:p03-w03-p03-007-two-cardinal-boundary}{A complete theory with a small infinite definable set}
Let \(\mathcal M\) be the structure with universe \(\omega_1\) and one unary
relation \(U^{\mathcal M}=\omega\), and let \(T_U\) be its complete theory.
For \(\theta=Uv_0\), the set \(\theta(\mathcal M)\) is infinite and of
smaller power than \(\mathcal M\). Consequently \(T_U\) is not
\(\aleph_1\)-categorical, and the conclusion of
\cref{thm:p03-w03-p03-007-definable-sets} fails without its categoricity
hypothesis.
\end{hbtheorem}

Vaught's results are stated for parameter-free \(\theta\), and the proof
of 5.4 uses this: the auxiliary theory names no elements, and the copying
step moves \(\theta\)-sets by arbitrary isomorphisms. Theorem~5.2 gives
existence only. Theorem~5.3 assumes that a prime model is given and does not
produce one. The cited theorem of Ehrenfeucht is used for the general
hypothesis of 5.2, and the companion use of 5.1 in Vaught's sixth section is
in \cref{sec:p03-vaught-never-two-presentation}.
